% ----------------------------------------------------------------------------
%                               Conclusão
% ----------------------------------------------------------------------------
\chapter{Conclusão}
\label{chap:conclusão}

O presente trabalho teve como objetivo central investigar e propor uma solução tecnológica voltada ao suporte pedagógico e ao manejo comportamental de estudantes com Transtorno do Déficit de Atenção com Hiperatividade (TDAH) e Transtorno Opositor Desafiador (TOD). Por meio de um levantamento rigoroso de requisitos e da subsequente codificação de ponta a ponta, foi possível materializar o sistema TEia, uma aplicação \textit{web} responsiva, implementada com arquitetura \textit{Mobile-First}, que visa atender às demandas latentes da comunidade escolar.

A principal contribuição deste projeto reside na proposição de uma transição do modelo de registro passivo para uma abordagem de manejo proativo. A interface desenvolvida, estruturada sob os princípios de usabilidade e baixa carga cognitiva, apresenta-se como uma alternativa viável para o dinâmico contexto educacional. A substituição de formulários textuais extensos por seleções ágeis (\textit{chips}) e a entrega imediata de estratégias de intervenção baseadas em evidências têm o potencial de reduzir significativamente a sobrecarga do corpo docente em momentos de crise.

Adicionalmente, a implementação de um perfil de acesso exclusivo para a Gestão Escolar, fundamentado em privilégios estritos de "somente leitura" e no rastreio dinâmico por meio de filtros combinados, consolidou o sistema TEia como uma solução de nível institucional. Ao ocultar as ferramentas de edição da coordenação, o sistema blindou a integridade dos prontuários e assegurou que a autoridade do registro comportamental permanecesse com o educador em sala de aula, respeitando rigorosamente as normativas do Estatuto da Criança e do Adolescente (ECA) e da Lei Geral de Proteção de Dados (LGPD).

Paralelamente, a estruturação de uma linha do tempo (\textit{timeline}) dedicada aos familiares democratiza o acesso à informação, promovendo uma comunicação transparente e livre de ruídos entre a escola e o ambiente doméstico. Essa arquitetura reafirma a importância da tecnologia não apenas como ferramenta de arquivamento, mas como ponte fundamental para a construção de uma rede de apoio colaborativa em prol do pleno desenvolvimento do estudante neurodivergente.

Visto que a implementação do código-fonte (utilizando \textit{Vanilla JS} e \textit{Firebase}) foi concluída com sucesso, e a validação do sistema funcional nas dependências do Instituto Federal de Brasília (IFB) foi iniciada por meio de um estudo de caso prático, as perspectivas do projeto evoluíram. Como trabalhos futuros, sugere-se a expansão da plataforma para abranger novas categorias de neurodivergência, a execução de testes em larga escala englobando múltiplos \textit{campi} da instituição e a evolução contínua das estratégias do assistente proativo a partir do \textit{feedback} longitudinal dos educadores e das famílias, visando consolidar o TEia como uma robusta tecnologia social na educação pública.